\section{Follow the history cursor through an overlap}
Start with output \texttt{41 42}, the ASCII bytes A and B. A match of length seven
and distance two reads from two bytes behind the growing output cursor. It appends
\texttt{41 42 41 42 41 42 41}; the complete result is \texttt{ABABABABA}. The third
copied byte comes from the first byte written by this same match. This dependence
explains why the forward copy is part of the format's meaning.

Keep three quantities separate: bytes produced so far, the match distance, and the
match length. With two bytes produced, distance two is available, even though a
length-seven match is longer than the current history. Capacity must accommodate all
seven new bytes. Distance three fails before the first copy load because its source
would precede the output start. Testing only the end destination cannot establish that
the history source exists.

The zlib header adds another distance bound. A stream advertising a 256-byte window
cannot request distance 257 even after producing thousands of bytes. A stream with
the largest permitted window can use distance 32768 once that much output exists.
The format's distance code and extra bits give the requested distance; the header
and current output position determine whether it is usable. Verify both bounds before
forming the source address. They answer different questions.

The dynamic header's length repeats have their own crossing rule. Literal/length and
distance code lengths form one combined sequence before they are separated into two
tables. For a declaration of 286 literal/length lengths and one distance length, a
repeat can fill literal/length slots 284 and 285 and then distance slot zero. Do not
reset the repeat count when the table kind changes. If the preceding length was zero,
operator 16 may repeat that zero just as it may repeat a nonzero length. Check that the
expanded run fits the combined count before storing any of its entries.
